% ==================== 1.5. MODELOS DE INTELIGENCIA ARTIFICIAL APLICADOS AL ANÁLISIS DE VOZ. ====================
        \section{Modelos de Inteligencia Artificial aplicados al análisis de voz}
            La aplicación de modelos de Inteligencia Artificial (IA) al análisis de la voz permite desarrollar sistemas capaces de detectar patrones acústicos y prosódicos asociados a enfermedades neurodegenerativas. Estos modelos se entrenan con características vocales previamente extraídas o directamente con representaciones espectrales de la señal, con el fin de clasificar a los individuos según su estado neurológico.
            En las investigaciones consultadas sobre el tema, los enfoques pueden agruparse en cuatro categorías principales: métodos clásicos de aprendizaje automático, modelos de aprendizaje profundo, modelos basados en embeddings o foundation models y enfoques de inteligencia artificial explicable (XAI).

            % ==================== 1.5.1. MÉTODOS CLÁSICOS DE APRENDIZAJE AUTOMÁTICO ====================
            \subsection{Métodos clásicos de aprendizaje automático}
                Los modelos clásicos de Machine Learning (ML) se han utilizado extensamente en la detección de Parkinson, Alzheimer y ELA debido a su facilidad de interpretabilidad, bajo coste computacional y eficacia en datasets con pocos datos. Entre los algoritmos más empleados destacan: Support Vector Machine (SVM), Random Forest (RF), K-Nearest Neighbors (KNN), Logistic Regression y AdaBoost, que operan sobre conjuntos de features previamente seleccionadas.

                En el ámbito del Parkinson, Tougui et al. (2020) emplea un sistema de clasificación usando SVM, KNN, RF y Extreme Gradient Boosting sobre 138 características en los dominios temporal, frecuencial y cepstral, logrando una precisión del 95,8\%, una sensibilidad del 95,3\% y una especificidad del 96,2\%.
                De forma parecida, Di Cesare et al. (2024) compararon los resultados de SVM, KNN y redes neuronales con características derivadas de MFCCs y Gammatone Cepstral Coefficients, logrando precisiones de hasta el 92,3\%, con sensibilidad y especificidad equilibradas (75–93\%).
                Otros estudios, como el de Kar y Campbell (2020), usando un modelo híbrido de regresión logística y red neuronal para detectar Parkinson en fases tempranas, a partir de medidas de Pitch Period Entropy y Spread1, obteniendo valores de sensibilidad del 95–96\% y precisión del 85–96\%.

                Para la enfermedad de Alzheimer, García-Gutiérrez et al. (2024) compararon varios algoritmos clásicos (RF, SVM, KNN y Extreme Gradient Boosting) en un conjunto de más de 1500 muestras de voz espontánea, alcanzando puntuaciones F1 entre 84\% y 92\%. 
                En la ELA, Tena et al. aplicaron SVM y Random Forest a componentes cuasi-periódicos de vocales sostenidas, obteniendo precisiones entre 78,7\% y 91,0\% y áreas bajo la curva (AUC) de 0,85–0,99, lo que demuestra la gran capacidad de estos métodos para detectar la afectación bulbar en su etapa inicial.

                En conjunto, los métodos clásicos siguen siendo una opción consolidada, sobretodo en los casos donde el tamaño de los datos es limitado y se dispone de features bien definidas. No obstante, su rendimiento suele verse afectado en datasets grandes o con características altamente correlacionadas.  En estos los modelos profundos ofrecen ventajas significativas.

            % ==================== 1.5.2. MODELOS DE APRENDIZAJE PROFUNDO ====================
            \subsection{Modelos de aprendizaje profundo}
                El Deep Learning ha transformado todo el campo del aprendizaje automático, incluido el del análisis de voz, permitiendo que los modelos aprendan representaciones directamente, a partir de la señal acústica, reduciendo la tediosa tarea de la extracción manual de características. Entre los enfoques más frecuentes destacan: las redes neuronales convolucionales (CNN), las redes recurrentes (RNN, LSTM, GRU) y los modelos híbridos CNN-RNN.

                En el ámbito del Parkinson, Worasawate et al. (2021) emplearon CNNs (LeNet-5, ResNet-50 y VGG-16), cuyas entradas eran los espectrogramas de vocales sostenidas /aa/ recogidos mediante smartphones. Con un dataset de 4051 muestras, alcanzaron precisiones del 97,7\% al 99,3\%, demostrando la capacidad de las CNN para capturar patrones espacio-temporales en señales ruidosas.
                Por su parte, Álvarez Marquina et al. (2022) aplicaron CNNs sobre medidas de formantes obtenidas de vocales sostenidas (/a/), obteniendo una precisión del 96\% y especificidad del 99\% en un conjunto de 1280 grabaciones, lo que confirma el potencial de las redes convolucionales para identificar alteraciones articulatorias.
                Luna-Ortiz et al. (2023), a través de una combinación de análisis dinámico no lineal supervisado y CNNs en dos bases de datos distintas, alcanzando resultados excepcionalmente altos (99,5–99,7\% de precisión).
                Finalmente, Muntaha et al. (2024) implementaron un modelo atencional híbrido LSTM-BiGRU con un componente de interpretabilidad (Explainable Boosting Machine), obteniendo precisiones entre 96,3\% y 97,5\%.

                Los modelos profundos también han mostrado resultados notables en Alzheimer y ELA. Moro-Velázquez et al. (2020, 2021) introdujeron técnicas de x-vectors e i-vectors, integradas en redes neuronales probabilísticas, alcanzando precisiones de 85–94\% y áreas bajo la curva (AUC) de 0,91–0,99 en Parkinson. Aunque su objetivo principal fue la diferenciación de fases clínicas de la enfermedad.

                En resumen, los modelos de aprendizaje profundo superan considerablemente a los métodos clásicos en datasets medianos o grandes, gracias a su capacidad para aprender representaciones jerárquicas complejas. Sin embargo, su rendimiento puede verse afectado en muestras pequeñas y requieren un control mayúsculo para evitar sobreajustes teniendo que realizar técnicas de regularización o validación cruzada.

            % ==================== 1.5.3. MODELOS BASADOS EN EMBEDDINGS Y FOUNDATION MODELS ====================
            \subsection{Modelos basados en embeddings y foundation models}
                En los últimos años, se ha observado una tendiencia hacia el uso de modelos auto-supervisados y de representación profunda del habla, capaces de aprender directamente de grandes volúmenes de datos no etiquetados. Estos modelos, denominados foundation models, generan embeddings que condensan información fonética, prosódica y semántica de manera más rica que las features tradicionales.

                En el estudio de Alzheimer, Agbavor y Liang (2022) usaron el modelo data2vec junto con clasificadores clásicos (SVM, RF, regresión logística), logrando áreas bajo la curva (AUC) de 0,835–0,846 y evidenciando que los embeddings auto-supervisados pueden sustituir las características manuales desempeñando un excelente papel.
                En el caso del Parkinson, Dao et al. (2025) exploraron wav2vec 2.0, Whisper y Seamless M4T para extraer representaciones de voz que fueron las entradas de distintos clasificadores. Dicho estudio reporta haber alcanzado un AUC del 91,4\%, confirmando la viabilidad de los foundation models para detectar etapas tempranas de la enfermedad.
               
                El uso de embeddings preentrenados ha demostrado conseguir mejorar la robustez frente a ruido, variabilidad entre hablantes y diferencias de idioma. Además, estos modelos reducen la necesidad de extracción de características y facilitan la transferencia entre tareas, abriendo nuevas posibilidades para la creación de sistemas universales de detección de enfermedades mediante voz.
            
            % ==================== 1.5.4. INTELIGENCIA ARTIFICIAL EXPLICABLE (XAI) Y MODELOS INTERPRETABLES ====================
            \subsection{Inteligencia Artificial explicable (XAI) y modelos interpretables}
                Un aspecto cada vez más relevante en la aplicación de IA a entornos clínicos es la interpretabilidad de los modelos. Los profesionales de la salud requieren entender qué características vocales influyen en la decisión del modelo para poder validar y confiar en los resultados.
                Shen et al. (2024) abordaron este desafío aplicando técnicas de Explainable AI a la detección temprana del Parkinson. Su enfoque combinó redes neuronales híbridas (CNN-RNN) con mecanismos de atención y visualización de relevancia, consiguiendo una precisión del 91,1\%, sensibilidad del 92,5\% y especificidad del 89,8\%, al tiempo que identificaban las features más influyentes (MFCCs, jitter, shimmer).
                De manera complementaria, Muntaha et al. (2024) integraron un componente de Explainable Boosting Machine (EBM) en su arquitectura LSTM-BiGRU, mostrando cómo la atención se concentraba en regiones de la señal asociadas a la disfonía y la falta de control respiratorio.

                Estos trabajos marcan una tendencia hacia modelos no solo precisos, sino también transparentes y clínicamente interpretables, condición necesaria para la futura integración de la IA en la práctica médica.

            % ==================== 1.5.5.  SINTESIS COMPARATIVA ====================
            \subsection{Sintesis Comparativa}
                El análisis conjunto de los doce estudios revisados permite identificar patrones claros:

                \begin{itemize}
                    \item Los \textbf{métodos clásicos} (SVM, RF, KNN) siguen siendo competitivos en datasets pequeños o con features bien seleccionadas, con precisiones medias del 80–95\%.
                    
                    \item Las \textbf{Redes profundas} (CNN, LSTM) ofrecen los mejores resultados globales (hasta 99\%) en escenarios con grandes volúmenes de datos y tareas controladas.
                    
                    \item Los \textbf{Modelos de embeddings y foundation models} representan la evolución natural del campo, destacando por su capacidad de generalización y menor dependencia del preprocesamiento manual.
                    
                    \item Los \textbf{enfoques explicables} emergen como una necesidad clave para la validación clínica y la adopción real en entornos sanitarios.

                \end{itemize}

                En conclusión, la IA aplicada al análisis de la voz alcanza niveles de rendimiento excepcionales en la detección de enfermedades neurodegenerativas, especialmente en Parkinson y Alzheimer. No obstante, siguen existiendo desafíos relacionados con la validación cruzada, la variabilidad de los datos y la interpretabilidad de los resultados. La integración de modelos fundacionales con técnicas de XAI se perfila como la línea más prometedora para futuras investigaciones y desarrollos clínicamente aplicables.
